\documentclass[../DoAn.tex]{subfiles}
\begin{document}
\section{Kết luận}
So với các công cụ đang được dùng để quản lý lịch phòng như bảng tính dùng chung hay lịch tài nguyên trực tuyến, RoomHub có ba điểm khác biệt. Thứ nhất, người dùng tìm phòng trực quan trên sơ đồ chữ U của từng tầng với màu trạng thái theo khung giờ, thay vì đọc danh sách lịch bận. Thứ hai, các quy tắc riêng của nhà trường như khoảng đặt trước, giới hạn số yêu cầu, lịch bảo trì, hạn hủy và hạn đổi phòng được kiểm tra tự động và có thể điều chỉnh ngay trong ứng dụng. Thứ ba, bước đặt phòng được nối liền với bước tiếp cận phòng: phòng khóa cơ hoặc thẻ từ có lịch hẹn và người nhận hộ, phòng khóa số có mã PIN tạm thời gửi thẳng tới khóa. Ngược lại, so với phần mềm đặt chỗ thương mại, RoomHub chưa có báo cáo thống kê, check-in bằng QR và tích hợp tài khoản trường.

Trong quá trình thực hiện, nhóm đã hoàn thành phân tích nghiệp vụ từ tài liệu quy trình phiên bản 1.1, thiết kế hướng đối tượng với các biểu đồ UML, xây dựng dịch vụ Java gồm 23 lớp và 12 điểm cuối REST, ứng dụng di động mười module với mười lăm màn hình, gateway khóa thông minh và bộ kiểm thử tự động. Ứng dụng đã được đóng gói thành APK Release và chạy độc lập trên Samsung M21. Ba đóng góp nổi bật là bộ quy tắc chống trùng lịch dựa trên phép giao khoảng thời gian kết hợp khóa bi quan, mô hình cấp quyền tạo mã theo cặp người dùng -- phòng tích hợp khóa DLWA12 qua MQTT, và kiến trúc hai chế độ dữ liệu dùng chung một bộ quy tắc.

Một số phần việc chưa hoàn thành gồm: bản phát hành hiện chạy ở chế độ cục bộ, dịch vụ Java chưa được triển khai trên máy chủ cố định và chưa có kiểm thử JUnit; các trạng thái Chờ bàn giao, Đang sử dụng, No-show cùng audit log đầy đủ trong tài liệu nghiệp vụ chưa được hiện thực; mật khẩu tài khoản và mã PIN còn lưu dạng rõ; hệ thống mới quản lý một khóa thông minh. Bài học rút ra của nhóm là cần chốt mô hình dữ liệu dùng chung giữa ứng dụng và máy chủ ngay từ đầu, viết quy tắc nghiệp vụ ở tầng dịch vụ có kiểm thử trước khi làm giao diện, và luôn kiểm chứng trên thiết bị thật vì nhiều lỗi chỉ xuất hiện khi chạy bản Release không có Metro.

\section{Hướng phát triển}
Trước hết, nhóm sẽ hoàn thiện các chức năng đã có. Dịch vụ Java sẽ được bật làm nguồn dữ liệu chính và triển khai lên máy chủ với PostgreSQL; khóa bi quan được chuyển xuống mức bản ghi phòng để loại bỏ hiện tượng phantom và tăng mức song song; bổ sung kiểm thử JUnit cùng bài kiểm thử 20 yêu cầu đồng thời. Về bảo mật, mật khẩu tài khoản sẽ được băm bằng BCrypt, phiên đăng nhập có thời hạn, mã PIN chỉ lưu dạng băm hoặc mã hóa, gateway bật xác minh chứng chỉ TLS và chỉ nhận lệnh từ máy chủ Java có xác thực.

Tiếp theo, hệ thống sẽ được mở rộng theo đầy đủ quy trình nghiệp vụ: bổ sung check-in, check-out bằng mã QR và tự động đánh dấu no-show sau thời gian chờ; quản lý số lượng, tình trạng từng thiết bị của phòng theo các quy tắc BR-13 đến BR-16 và cảnh báo booking bị ảnh hưởng khi thiết bị hỏng; ghi audit log cho mọi thao tác bàn giao, cấp mã, hủy, đổi phòng; quản lý thẻ từ với vòng đời tại tủ, giữ chỗ, bàn giao, trả và đối soát. Ở mức trải nghiệm, nhóm dự kiến bổ sung gợi ý phòng thay thế khi phòng mong muốn đã kín, chỉ đường từ thang máy tới phòng trên sơ đồ, thống kê mức độ sử dụng phòng cho đơn vị quản lý, hỗ trợ nhiều khóa thông minh và tích hợp đăng nhập bằng tài khoản trường.
\end{document}
